\documentclass[11pt,a4paper]{article}
\usepackage[utf8]{inputenc}
\usepackage[T1]{fontenc}
\usepackage[spanish,es-minimal,es-noshorthands]{babel}
\AtBeginDocument{\renewcommand{\sin}{\operatorname{sen}}\renewcommand{\arcsin}{\operatorname{arcsen}}}
\usepackage{lmodern}
\usepackage[sfdefault,lining]{FiraSans}
\usepackage[a4paper,top=2.3cm,bottom=2.2cm,left=1.8cm,right=1.8cm,headheight=1.3cm,headsep=0.4cm,footskip=1cm]{geometry}
\usepackage{amsmath,amssymb}
\usepackage{icomma}
\usepackage{graphicx}
\usepackage[table]{xcolor}
\usepackage{tikz}
\usetikzlibrary{calc}
\usepackage{fontawesome5}
\usepackage{ccicons}
\usepackage{fancyhdr}
\usepackage{lastpage}
\usepackage{tasks}
\usepackage{needspace}
\usepackage[most]{tcolorbox}
\usepackage{array}
\usepackage[hidelinks]{hyperref}
\AtBeginDocument{\hypersetup{pdftitle={Ficha de ejercicios de la unidad 3: Sistemas de ecuaciones e inecuaciones. 1º Bach.},pdfsubject={Matemáticas CCSS 1º Bach.},pdfkeywords={sistemas de ecuaciones lineales, método de Gauss, sistemas no lineales, inecuaciones, primer grado, segundo grado, sistemas de inecuaciones, semiplanos, regiones del plano, programación lineal}}}

\definecolor{navy}{HTML}{1B2748}
\definecolor{accent}{HTML}{E8603A}
\definecolor{teal}{HTML}{1F9E90}
\definecolor{soft}{HTML}{EEF1F8}
\definecolor{muted}{HTML}{5D6685}

\setlength{\parindent}{0pt}
\setlength{\parskip}{3pt}
\settasks{label=\alph*), label-format=\bfseries\color{navy}, label-width=1.3em, item-indent=1.8em, column-sep=0.6em, after-item-skip=5pt, before-skip=3pt, after-skip=4pt}

% ---------- capçalera i peu ----------
\newcommand{\llicencia}{\href{https://creativecommons.org/licenses/by-nc-sa/4.0/deed.es}{\ccbyncsa}\ \ Licencia Creative Commons Atribución-NoComercial-CompartirIgual 4.0 Internacional (CC BY-NC-SA 4.0)}
\pagestyle{fancy}
\fancyhf{}
\fancyhead[L]{\raisebox{-0.35\height}{\includegraphics[height=0.95cm]{logo-gva.png}}}
\fancyhead[C]{\small\color{muted}Ficha de ejercicios · Unidad 3: Sistemas de ecuaciones e inecuaciones · 1\textsuperscript{o} Bach.}
\fancyhead[R]{\raisebox{-0.35\height}{\includegraphics[height=0.95cm]{logo-vermellar.png}}}
\fancyfoot[L]{\scriptsize\color{muted}\llicencia}
\fancyfoot[R]{\small\color{muted}\thepage}
\renewcommand{\headrulewidth}{0.4pt}
\renewcommand{\headrule}{\color{muted!40}\hrule height \headrulewidth}
\fancypagestyle{primera}{\fancyhf{}\renewcommand{\headrulewidth}{0pt}
  \fancyfoot[L]{\scriptsize\color{muted}\llicencia}
  \fancyfoot[R]{\small\color{muted}\thepage}}

% ---------- nivells i exercicis ----------
\newcommand{\estrelles}[1]{{\color{accent}\ifcase#1\or\faStar\faStar[regular]\faStar[regular]\or\faStar\faStar\faStar[regular]\or\faStar\faStar\faStar\fi}}
\newcounter{exer}
\newcommand{\ex}[1]{\par\needspace{4\baselineskip}\medskip\refstepcounter{exer}\noindent{\color{navy}\bfseries\theexer.}\ \ }

\newcommand{\seccio}[2]{\par\needspace{16\baselineskip}\bigskip
  \begin{tcolorbox}[colback=navy,colframe=navy,coltext=white,arc=2mm,boxrule=0pt,left=3mm,right=3mm,top=1.2mm,bottom=1.2mm]
  {\large\bfseries #1\quad #2}
  \end{tcolorbox}\vspace{-2pt}}

\newcommand{\exf}[1]{\par\medskip\refstepcounter{exer}\noindent{\color{navy}\bfseries\theexer.}\ \ }
\providecommand{\lt}{<}\providecommand{\gt}{>}
\newcolumntype{C}[1]{>{\centering\arraybackslash}p{#1}}
% quadre de recordatori (només a la fitxa d'exercicis)
\newcommand{\recorda}[1]{\par\smallskip\begin{tcolorbox}[colback=accent!7,colframe=accent!55,arc=2mm,boxrule=0.5pt,left=2.5mm,right=2.5mm,top=1mm,bottom=1mm]\small{\color{accent}\faLightbulb[regular]}\ \textbf{Recuerda:}\ #1\end{tcolorbox}}
\newcommand{\res}[1]{#1}\newcommand{\rest}[1]{#1}

\def\fmneed{11}% que el títol d'una secció no quede sol al final del full
% =====================================================================
% Capçalera FentMates de les fitxes: un color per a cada curs.
% Els colors són els de la pàgina inicial (LEVELS a assets/js/materials.js);
% el logo ve de _eines/logo/pdf/ (python3 _eines/logo/genera.py).
%
% Ús, al final del preàmbul de la unitat:
%   % =====================================================================
% Capçalera FentMates de les fitxes: un color per a cada curs.
% Els colors són els de la pàgina inicial (LEVELS a assets/js/materials.js);
% el logo ve de _eines/logo/pdf/ (python3 _eines/logo/genera.py).
%
% Ús, al final del preàmbul de la unitat:
%   % =====================================================================
% Capçalera FentMates de les fitxes: un color per a cada curs.
% Els colors són els de la pàgina inicial (LEVELS a assets/js/materials.js);
% el logo ve de _eines/logo/pdf/ (python3 _eines/logo/genera.py).
%
% Ús, al final del preàmbul de la unitat:
%   \input{../comu/fentmates.tex}
%   \fmcurs{1eso}
% i, al començament del document:
%   \capcalera{Matemàtiques}{1r ESO}{Fitxa d'exercicis · Unitat 2}{Divisibilitat}
% Opcional abans de \input: \def\fmneed{14} (espai mínim abans de cada secció).
% =====================================================================
\makeatletter
\def\fm@logodir{../../logo/pdf}
\providecommand{\fmneed}{10}
\def\fm@cursnom{}

% clau del curs, color i color del text damunt del color
\newcommand{\fm@curs}[3]{\@namedef{fm@color@#1}{#2}\@namedef{fm@text@#1}{#3}}
\fm@curs{1eso}{E6B84A}{navy}
\fm@curs{2eso}{3B82F6}{white}
\fm@curs{3eso}{FF7A50}{white}
\fm@curs{4eso-a}{A78BFA}{white}
\fm@curs{4eso-b}{6D28D9}{white}
\fm@curs{1batx-mi}{9B1C1C}{white}
\fm@curs{1batx-ccss}{EF6B6B}{white}
\fm@curs{2batx-mii}{1E6B47}{white}
\fm@curs{2batx-ccss}{2BB3A3}{white}

\newcommand{\fmcurs}[1]{%
  \@ifundefined{fm@color@#1}{\PackageError{fentmates}{Curs desconegut: #1}{}}{}%
  \edef\fm@tmp{\noexpand\definecolor{curs}{HTML}{\@nameuse{fm@color@#1}}%
               \noexpand\colorlet{curstext}{\@nameuse{fm@text@#1}}}\fm@tmp
  \colorlet{cursfosc}{curs!72!black}%
  \def\fm@logo{\fm@logodir/vertical-#1.pdf}%
  \def\fm@marca{\fm@logodir/marca-#1.pdf}}

% ---------- capçalera de la primera pàgina ----------
% #1 matèria (petit), #2 curs, #3 tipus de document i unitat, #4 títol
\newcommand{\capcalera}[4]{%
  \gdef\fm@cursnom{#2}%
  \par\noindent
  \begin{tcolorbox}[colback=curs,colframe=curs,boxrule=0pt,arc=2.2mm,width=3.5cm,height=2.15cm,
      valign=center,left=3mm,right=2mm,top=1mm,bottom=1mm,nobeforeafter,box align=center]
    \color{curstext}{\footnotesize #1}\par\vspace{1pt}{\fontsize{19}{21}\selectfont\bfseries #2\par}%
  \end{tcolorbox}\hspace{5mm}%
  \begin{minipage}[c]{\dimexpr\linewidth-3.5cm-5mm-2.9cm\relax}
    {\color{cursfosc}\small\bfseries\MakeUppercase{#3}}\par\vspace{2pt}
    {\color{navy}\fontsize{20}{23}\selectfont\bfseries #4\par}
  \end{minipage}\hfill
  \begin{minipage}[c]{2.6cm}\raggedleft\includegraphics[width=2.6cm]{\fm@logo}\end{minipage}\par
  \vspace{2pt}}

% ---------- barres de secció del color del curs ----------
\renewcommand{\seccio}[2]{\par\needspace{\fmneed\baselineskip}\bigskip
  \begin{tcolorbox}[colback=curs,colframe=curs,coltext=curstext,arc=2mm,boxrule=0pt,left=3mm,right=3mm,top=1.2mm,bottom=1.2mm]
  {\large\bfseries #1\quad #2}
  \end{tcolorbox}\vspace{-2pt}}

% ---------- estrelles de dificultat del color del curs (a Batxillerat no se'n posen) ----------
\@ifundefined{estrelles}{}{%
  \renewcommand{\estrelles}[1]{{\color{curs}\ifcase#1\or\faStar\faStar[regular]\faStar[regular]\or\faStar\faStar\faStar[regular]\or\faStar\faStar\faStar\fi}}}

% ---------- capçalera i peu de les pàgines ----------
% primera pàgina: logos institucionals al peu
\fancypagestyle{primera}{\fancyhf{}\renewcommand{\headrulewidth}{0pt}
  \fancyfoot[L]{\raisebox{-0.32\height}{\includegraphics[height=0.78cm]{logo-gva.png}}\hspace{5mm}%
    \raisebox{-0.32\height}{\includegraphics[height=0.78cm]{logo-vermellar.png}}}
  \fancyfoot[R]{\scriptsize\color{muted}\href{https://creativecommons.org/licenses/by-nc-sa/4.0/deed.ca}{\ccbyncsa}\ \ CC BY-NC-SA 4.0 · fentmates.org\qquad\small\thepage}}
% la resta: el símbol a l'esquerra i el curs a la dreta, amb el títol al centre com abans
\fancyhead[L]{\raisebox{-0.3\height}{\includegraphics[height=0.62cm]{\fm@marca}}}
\fancyhead[R]{\small\bfseries\color{cursfosc}\fm@cursnom}
\renewcommand{\headrule}{\color{curs!45}\hrule height \headrulewidth}
\makeatother

%   \fmcurs{1eso}
% i, al començament del document:
%   \capcalera{Matemàtiques}{1r ESO}{Fitxa d'exercicis · Unitat 2}{Divisibilitat}
% Opcional abans de \input: \def\fmneed{14} (espai mínim abans de cada secció).
% =====================================================================
\makeatletter
\def\fm@logodir{../../logo/pdf}
\providecommand{\fmneed}{10}
\def\fm@cursnom{}

% clau del curs, color i color del text damunt del color
\newcommand{\fm@curs}[3]{\@namedef{fm@color@#1}{#2}\@namedef{fm@text@#1}{#3}}
\fm@curs{1eso}{E6B84A}{navy}
\fm@curs{2eso}{3B82F6}{white}
\fm@curs{3eso}{FF7A50}{white}
\fm@curs{4eso-a}{A78BFA}{white}
\fm@curs{4eso-b}{6D28D9}{white}
\fm@curs{1batx-mi}{9B1C1C}{white}
\fm@curs{1batx-ccss}{EF6B6B}{white}
\fm@curs{2batx-mii}{1E6B47}{white}
\fm@curs{2batx-ccss}{2BB3A3}{white}

\newcommand{\fmcurs}[1]{%
  \@ifundefined{fm@color@#1}{\PackageError{fentmates}{Curs desconegut: #1}{}}{}%
  \edef\fm@tmp{\noexpand\definecolor{curs}{HTML}{\@nameuse{fm@color@#1}}%
               \noexpand\colorlet{curstext}{\@nameuse{fm@text@#1}}}\fm@tmp
  \colorlet{cursfosc}{curs!72!black}%
  \def\fm@logo{\fm@logodir/vertical-#1.pdf}%
  \def\fm@marca{\fm@logodir/marca-#1.pdf}}

% ---------- capçalera de la primera pàgina ----------
% #1 matèria (petit), #2 curs, #3 tipus de document i unitat, #4 títol
\newcommand{\capcalera}[4]{%
  \gdef\fm@cursnom{#2}%
  \par\noindent
  \begin{tcolorbox}[colback=curs,colframe=curs,boxrule=0pt,arc=2.2mm,width=3.5cm,height=2.15cm,
      valign=center,left=3mm,right=2mm,top=1mm,bottom=1mm,nobeforeafter,box align=center]
    \color{curstext}{\footnotesize #1}\par\vspace{1pt}{\fontsize{19}{21}\selectfont\bfseries #2\par}%
  \end{tcolorbox}\hspace{5mm}%
  \begin{minipage}[c]{\dimexpr\linewidth-3.5cm-5mm-2.9cm\relax}
    {\color{cursfosc}\small\bfseries\MakeUppercase{#3}}\par\vspace{2pt}
    {\color{navy}\fontsize{20}{23}\selectfont\bfseries #4\par}
  \end{minipage}\hfill
  \begin{minipage}[c]{2.6cm}\raggedleft\includegraphics[width=2.6cm]{\fm@logo}\end{minipage}\par
  \vspace{2pt}}

% ---------- barres de secció del color del curs ----------
\renewcommand{\seccio}[2]{\par\needspace{\fmneed\baselineskip}\bigskip
  \begin{tcolorbox}[colback=curs,colframe=curs,coltext=curstext,arc=2mm,boxrule=0pt,left=3mm,right=3mm,top=1.2mm,bottom=1.2mm]
  {\large\bfseries #1\quad #2}
  \end{tcolorbox}\vspace{-2pt}}

% ---------- estrelles de dificultat del color del curs (a Batxillerat no se'n posen) ----------
\@ifundefined{estrelles}{}{%
  \renewcommand{\estrelles}[1]{{\color{curs}\ifcase#1\or\faStar\faStar[regular]\faStar[regular]\or\faStar\faStar\faStar[regular]\or\faStar\faStar\faStar\fi}}}

% ---------- capçalera i peu de les pàgines ----------
% primera pàgina: logos institucionals al peu
\fancypagestyle{primera}{\fancyhf{}\renewcommand{\headrulewidth}{0pt}
  \fancyfoot[L]{\raisebox{-0.32\height}{\includegraphics[height=0.78cm]{logo-gva.png}}\hspace{5mm}%
    \raisebox{-0.32\height}{\includegraphics[height=0.78cm]{logo-vermellar.png}}}
  \fancyfoot[R]{\scriptsize\color{muted}\href{https://creativecommons.org/licenses/by-nc-sa/4.0/deed.ca}{\ccbyncsa}\ \ CC BY-NC-SA 4.0 · fentmates.org\qquad\small\thepage}}
% la resta: el símbol a l'esquerra i el curs a la dreta, amb el títol al centre com abans
\fancyhead[L]{\raisebox{-0.3\height}{\includegraphics[height=0.62cm]{\fm@marca}}}
\fancyhead[R]{\small\bfseries\color{cursfosc}\fm@cursnom}
\renewcommand{\headrule}{\color{curs!45}\hrule height \headrulewidth}
\makeatother

%   \fmcurs{1eso}
% i, al començament del document:
%   \capcalera{Matemàtiques}{1r ESO}{Fitxa d'exercicis · Unitat 2}{Divisibilitat}
% Opcional abans de \input: \def\fmneed{14} (espai mínim abans de cada secció).
% =====================================================================
\makeatletter
\def\fm@logodir{../../logo/pdf}
\providecommand{\fmneed}{10}
\def\fm@cursnom{}

% clau del curs, color i color del text damunt del color
\newcommand{\fm@curs}[3]{\@namedef{fm@color@#1}{#2}\@namedef{fm@text@#1}{#3}}
\fm@curs{1eso}{E6B84A}{navy}
\fm@curs{2eso}{3B82F6}{white}
\fm@curs{3eso}{FF7A50}{white}
\fm@curs{4eso-a}{A78BFA}{white}
\fm@curs{4eso-b}{6D28D9}{white}
\fm@curs{1batx-mi}{9B1C1C}{white}
\fm@curs{1batx-ccss}{EF6B6B}{white}
\fm@curs{2batx-mii}{1E6B47}{white}
\fm@curs{2batx-ccss}{2BB3A3}{white}

\newcommand{\fmcurs}[1]{%
  \@ifundefined{fm@color@#1}{\PackageError{fentmates}{Curs desconegut: #1}{}}{}%
  \edef\fm@tmp{\noexpand\definecolor{curs}{HTML}{\@nameuse{fm@color@#1}}%
               \noexpand\colorlet{curstext}{\@nameuse{fm@text@#1}}}\fm@tmp
  \colorlet{cursfosc}{curs!72!black}%
  \def\fm@logo{\fm@logodir/vertical-#1.pdf}%
  \def\fm@marca{\fm@logodir/marca-#1.pdf}}

% ---------- capçalera de la primera pàgina ----------
% #1 matèria (petit), #2 curs, #3 tipus de document i unitat, #4 títol
\newcommand{\capcalera}[4]{%
  \gdef\fm@cursnom{#2}%
  \par\noindent
  \begin{tcolorbox}[colback=curs,colframe=curs,boxrule=0pt,arc=2.2mm,width=3.5cm,height=2.15cm,
      valign=center,left=3mm,right=2mm,top=1mm,bottom=1mm,nobeforeafter,box align=center]
    \color{curstext}{\footnotesize #1}\par\vspace{1pt}{\fontsize{19}{21}\selectfont\bfseries #2\par}%
  \end{tcolorbox}\hspace{5mm}%
  \begin{minipage}[c]{\dimexpr\linewidth-3.5cm-5mm-2.9cm\relax}
    {\color{cursfosc}\small\bfseries\MakeUppercase{#3}}\par\vspace{2pt}
    {\color{navy}\fontsize{20}{23}\selectfont\bfseries #4\par}
  \end{minipage}\hfill
  \begin{minipage}[c]{2.6cm}\raggedleft\includegraphics[width=2.6cm]{\fm@logo}\end{minipage}\par
  \vspace{2pt}}

% ---------- barres de secció del color del curs ----------
\renewcommand{\seccio}[2]{\par\needspace{\fmneed\baselineskip}\bigskip
  \begin{tcolorbox}[colback=curs,colframe=curs,coltext=curstext,arc=2mm,boxrule=0pt,left=3mm,right=3mm,top=1.2mm,bottom=1.2mm]
  {\large\bfseries #1\quad #2}
  \end{tcolorbox}\vspace{-2pt}}

% ---------- estrelles de dificultat del color del curs (a Batxillerat no se'n posen) ----------
\@ifundefined{estrelles}{}{%
  \renewcommand{\estrelles}[1]{{\color{curs}\ifcase#1\or\faStar\faStar[regular]\faStar[regular]\or\faStar\faStar\faStar[regular]\or\faStar\faStar\faStar\fi}}}

% ---------- capçalera i peu de les pàgines ----------
% primera pàgina: logos institucionals al peu
\fancypagestyle{primera}{\fancyhf{}\renewcommand{\headrulewidth}{0pt}
  \fancyfoot[L]{\raisebox{-0.32\height}{\includegraphics[height=0.78cm]{logo-gva.png}}\hspace{5mm}%
    \raisebox{-0.32\height}{\includegraphics[height=0.78cm]{logo-vermellar.png}}}
  \fancyfoot[R]{\scriptsize\color{muted}\href{https://creativecommons.org/licenses/by-nc-sa/4.0/deed.ca}{\ccbyncsa}\ \ CC BY-NC-SA 4.0 · fentmates.org\qquad\small\thepage}}
% la resta: el símbol a l'esquerra i el curs a la dreta, amb el títol al centre com abans
\fancyhead[L]{\raisebox{-0.3\height}{\includegraphics[height=0.62cm]{\fm@marca}}}
\fancyhead[R]{\small\bfseries\color{cursfosc}\fm@cursnom}
\renewcommand{\headrule}{\color{curs!45}\hrule height \headrulewidth}
\makeatother

\fmcurs{1batx-ccss}

\begin{document}
\thispagestyle{primera}

\capcalera{Matemáticas CCSS}{1\textsuperscript{o} Bach.}{Ficha de ejercicios · Unidad 3}{\raggedright Sistemas de ecuaciones e inecuaciones}
\vspace{10pt}
Nombre y apellidos: \rule{6.9cm}{0.4pt}\hfill Grupo: \rule{1.6cm}{0.4pt}\hfill Fecha: \rule{2.2cm}{0.4pt}
\vspace{8pt}

\begin{tcolorbox}[colback=soft,colframe=soft,arc=2mm,left=3mm,right=3mm,top=1.5mm,bottom=1.5mm]
\small Resuelve los ejercicios en el cuaderno mostrando todos los pasos. Escribe las soluciones de las inecuaciones en forma de intervalo y represéntalas en la recta real; en las de dos incógnitas, dibuja la región solución. En los problemas, define las incógnitas y responde con una frase.
\end{tcolorbox}

% =====================================================
\seccio{1}{Sistemas de ecuaciones lineales}

\exf{1} Resuelve por el método que prefieras (sustitución, igualación o reducción):\begin{tasks}(4)
  \task $\left\{\begin{array}{l}2x+y=7\\x-y=2\end{array}\right.$
  \task $\left\{\begin{array}{l}3x-2y=4\\x+y=3\end{array}\right.$
  \task $\left\{\begin{array}{l}x+3y=10\\2x-y=-1\end{array}\right.$
  \task $\left\{\begin{array}{l}4x+3y=1\\2x-y=3\end{array}\right.$
\end{tasks}

\ex{2} Quita paréntesis, simplifica cada ecuación hasta la forma $ax+by=c$ y resuelve:\begin{tasks}(3)
  \task $\left\{\begin{array}{l}2(x+1)-3y=7\\x-(y-2)=5\end{array}\right.$
  \task $\left\{\begin{array}{l}5(x-2)+y=3(y+1)\\2x-3(y-1)=6\end{array}\right.$
  \task $\left\{\begin{array}{l}4(x+y)-3(x-y)=17\\2(x-3y)+y=-4\end{array}\right.$
\end{tasks}\recorda{Cuidado con el signo $-$ delante de un paréntesis: cambia el signo de todos los términos de dentro.}

\ex{2} Quita denominadores y resuelve:\begin{tasks}(3)
  \task $\left\{\begin{array}{l}\dfrac x3+\dfrac y2=3\\[5pt]\dfrac{x-y}{2}=2\end{array}\right.$
  \task $\left\{\begin{array}{l}\dfrac{x+1}{2}-\dfrac{y-1}{3}=2\\[5pt]\dfrac{x+y}{3}=3\end{array}\right.$
  \task $\left\{\begin{array}{l}\dfrac x2-\dfrac y5=1\\[5pt]\dfrac{x+2}{3}+\dfrac y5=3\end{array}\right.$
\end{tasks}\recorda{Multiplica los dos miembros de cada ecuación por el m.c.m. de sus denominadores.}

\ex{2} Resuelve y clasifica cada sistema. Interprétalo gráficamente:\begin{tasks}(3)
  \task $\left\{\begin{array}{l}\dfrac{x+1}{2}-y=1\\[5pt]3x-6y=3\end{array}\right.$
  \task $\left\{\begin{array}{l}2(x-y)=3\\x-y=1\end{array}\right.$
  \task $\left\{\begin{array}{l}x+2y=4\\2x+y=5\end{array}\right.$
\end{tasks}

\ex{2} En una granja hay gallinas y conejos. En total hay 30 cabezas y 80 patas. ¿Cuántas gallinas y cuántos conejos hay?

\ex{2} Un cine vende entradas normales a 7 € y reducidas a 5 €. Un día ha vendido 120 entradas y ha recaudado 720 €. ¿Cuántas ha vendido de cada tipo?


% =====================================================
\seccio{2}{Método de Gauss}

\exf{1} Resuelve estos sistemas escalonados, de abajo hacia arriba:\begin{tasks}(2)
  \task $\left\{\begin{array}{l}x+y+z=6\\y-z=-1\\2z=6\end{array}\right.$
  \task $\left\{\begin{array}{l}2x-y+z=3\\3y-z=1\\z=2\end{array}\right.$
\end{tasks}

\ex{2} Resuelve por el método de Gauss y clasifica:\begin{tasks}(2)
  \task $\left\{\begin{array}{l}x+2y-3z=7\\2x+y-z=6\\3x-y-z=6\end{array}\right.$
  \task $\left\{\begin{array}{l}x-y+z=2\\2x+y-z=1\\x+2y+3z=14\end{array}\right.$
  \task $\left\{\begin{array}{l}x+y-z=5\\3x-2y-2z=5\\2x-3y-z=3\end{array}\right.$
  \task $\left\{\begin{array}{l}2x+y+z=5\\2x-y+3z=3\\3x+2y+z=8\end{array}\right.$
\end{tasks}\recorda{Haz ceros debajo de la diagonal. La última fila $(0\ \ 0\ \ m\,|\,n)$ decide: $m\neq0$ SCD; $m=n=0$ SCI; $m=0,\ n\neq0$ SI.}

\ex{2} En una papelería, 2 bolígrafos, 1 libreta y 1 carpeta cuestan 9 €; 1 bolígrafo, 2 libretas y 1 carpeta, 10 €, y 1 bolígrafo, 1 libreta y 2 carpetas, 13 €. ¿Cuánto cuesta cada artículo?

\ex{3} Un examen tipo test tiene 30 preguntas. Cada acierto suma 2 puntos, cada error resta 1 punto y las preguntas en blanco no cuentan. Una alumna ha sacado 36 puntos y ha tenido el doble de aciertos que de errores y preguntas en blanco juntas. ¿Cuántas preguntas ha acertado, ha fallado y ha dejado en blanco?


% =====================================================
\seccio{3}{Sistemas de ecuaciones no lineales}

\exf{2} Resuelve por el método de sustitución:\begin{tasks}(4)
  \task $\left\{\begin{array}{l}x+y=7\\xy=12\end{array}\right.$
  \task $\left\{\begin{array}{l}y=x^2-4\\y=x+2\end{array}\right.$
  \task $\left\{\begin{array}{l}2x-y=1\\x^2+y^2=2\end{array}\right.$
  \task $\left\{\begin{array}{l}x+y=2\\xy=5\end{array}\right.$
\end{tasks}\recorda{Despeja una incógnita de la ecuación lineal y sustitúyela en la otra. Por cada valor que encuentres, calcula el de la otra incógnita.}

\ex{2} Resuelve por el método de reducción:\begin{tasks}(2)
  \task $\left\{\begin{array}{l}x^2+y^2=13\\x^2-y^2=5\end{array}\right.$
  \task $\left\{\begin{array}{l}x^2+2y^2=18\\2x^2-y^2=31\end{array}\right.$
\end{tasks}

\ex{2} El perímetro de un rectángulo mide 34 cm y su diagonal, 13 cm. ¿Cuánto miden los lados?

\ex{2} La diferencia de dos números es 3 y su producto, 40. ¿Cuáles son?


% =====================================================
\seccio{4}{Inecuaciones de primer grado}

\exf{1} Resuelve las inecuaciones siguientes y representa su solución:\begin{tasks}(3)
  \task $3x-7\gt5$
  \task $2x+5\le4x-1$
  \task $5-2x\ge11$
  \task $4(x-1)\lt2x+6$
  \task $3(2-x)\ge2(x+8)$
  \task $x-2(x+3)\gt4-3x$
\end{tasks}\recorda{Si multiplicas o divides por un número negativo, la desigualdad cambia de sentido.}

\ex{2} Resuelve las inecuaciones con denominadores:\begin{tasks}(2)
  \task $\dfrac{x}{2}+1\le\dfrac{x}{3}+2$
  \task $\dfrac{x-1}{3}-\dfrac{x+2}{4}\gt1$
  \task $\dfrac{2x+1}{5}\ge\dfrac{x-3}{2}$
  \task $x-\dfrac{x+1}{2}\lt\dfrac{x-2}{3}$
\end{tasks}

\ex{2} Resuelve e interpreta el resultado:\begin{tasks}(2)
  \task $2(x+3)\gt2x+1$
  \task $3x-1\le3(x-2)$
\end{tasks}

\ex{2} Un gimnasio A cobra 30 € de matrícula y 25 € al mes, y un gimnasio B no cobra matrícula pero cuesta 30 € al mes. ¿A partir de cuántos meses es más barato el gimnasio A?

\ex{2} Un lado de un rectángulo mide 8 cm. ¿Cuánto puede medir el otro lado para que el perímetro no supere los 30 cm?

\ex{2} Un vendedor cobra 900 € fijos al mes más un 5\,\% de las ventas. ¿Cuánto tiene que vender para cobrar como mínimo 1\,400 €?


% =====================================================
\seccio{5}{Inecuaciones de segundo grado}

\exf{1} Resuelve:\begin{tasks}(3)
  \task $x^2-3x-4\lt0$
  \task $x^2-x-6\ge0$
  \task $x^2-16\gt0$
  \task $x^2+3x\le0$
  \task $-x^2+4x-3\gt0$
  \task $2x^2-5x+2\le0$
\end{tasks}\recorda{Calcula las raíces, divide la recta en intervalos y estudia el signo en cada uno (o piensa en la parábola).}

\ex{2} Resuelve estos casos especiales:\begin{tasks}(2)
  \task $x^2+x+1\gt0$
  \task $x^2+2x+4\le0$
  \task $x^2-4x+4\gt0$
  \task $x^2-6x+9\le0$
\end{tasks}

\ex{2} Pásalo todo a un miembro y resuelve:\begin{tasks}(3)
  \task $x(x-2)\gt3$
  \task $(x+1)^2\le2x+5$
  \task $\dfrac{x^2}{2}-x\ge4$
\end{tasks}

\ex{2} El beneficio de una empresa viene dado por \mbox{$B(x)=-x^2+10x-16$}, en miles de euros, donde $x$ son los centenares de unidades que vende. ¿Cuántas unidades tiene que vender para tener beneficios ($B(x)\gt0$)?

\ex{2} La suma de un número y su cuadrado es menor que 12. ¿Qué valores puede tomar el número?


% =====================================================
\seccio{6}{Inecuaciones de grado superior}

\exf{2} Factoriza y resuelve con una tabla de signos:\begin{tasks}(2)
  \task $x^3-4x\ge0$
  \task $(x-1)(x+2)(x-3)\lt0$
  \task $x^3-x^2-2x\gt0$
  \task $x^3-6x^2+11x-6\ge0$
\end{tasks}

\ex{3} Cuidado con las raíces dobles: un factor al cuadrado no cambia de signo. Resuelve:\begin{tasks}(3)
  \task $x^3+x^2-x-1\le0$
  \task $x^2(x-3)\le0$
  \task $(x+2)^2(x-1)\gt0$
\end{tasks}

\ex{3} Resuelve la inecuación bicuadrada $x^4-5x^2+4\gt0$.


% =====================================================
\seccio{7}{Sistemas de inecuaciones con una incógnita}

\exf{2} Resuelve los sistemas de inecuaciones:\begin{tasks}(3)
  \task $\left\{\begin{array}{l}x+3\gt0\\2x-1\le5\end{array}\right.$
  \task $\left\{\begin{array}{l}3x-2\lt x+4\\5-x\le2x+2\end{array}\right.$
  \task $\left\{\begin{array}{l}2x+1\gt x+4\\3x-5\lt x+1\end{array}\right.$
  \task $\left\{\begin{array}{l}\dfrac{x+1}{2}\ge x-1\\4x+3\gt x\end{array}\right.$
  \task $\left\{\begin{array}{l}x^2-4\lt0\\2x+1\ge0\end{array}\right.$
  \task $\left\{\begin{array}{l}x^2-5x+4\le0\\x^2-9\gt0\end{array}\right.$
\end{tasks}\recorda{La solución del sistema es la intersección de las soluciones: dibújalas en rectas una sobre la otra.}

\ex{2} Halla los números naturales tales que su triple menos 5 es mayor que 10 y su doble más 1 es menor que 21.

\ex{2} Una familia quiere recorrer entre 400 y 600 km en un día a una velocidad media de 80 km/h. ¿Cuántas horas tiene que conducir?


% =====================================================
\seccio{8}{Inecuaciones con dos incógnitas}

\exf{1} Representa la solución de las inecuaciones:\begin{tasks}(4)
  \task $x+y\le4$
  \task $2x-y\gt2$
  \task $y\lt3x$
  \task $x\ge-2$
\end{tasks}\recorda{Dibuja la recta (continua si es $\le,\ge$; discontinua si es $\lt,\gt$) y prueba un punto que no esté en ella, como $O(0,0)$.}

\ex{1} Comprueba si los puntos $A(1,2)$, $B(3,1)$ y $C(-1,4)$ son solución de la inecuación $2x+y\le5$.

\ex{2} Con 20 € como máximo queremos comprar bolígrafos de 2 € y cuadernos de 4 €. Escribe las inecuaciones que tienen que cumplir el número de bolígrafos, $x$, y de cuadernos, $y$, y represéntalas. ¿Podemos comprar 4 bolígrafos y 3 cuadernos?


% =====================================================
\seccio{9}{Sistemas de inecuaciones con dos incógnitas}

\exf{2} Resuelve los sistemas, halla los vértices de la región solución y di si es acotada:\begin{tasks}(3)
  \task $\left\{\begin{array}{l}x+y\le4\\x\ge0,\ y\ge0\end{array}\right.$
  \task $\left\{\begin{array}{l}x+2y\le8\\x\le4\\x\ge0,\ y\ge0\end{array}\right.$
  \task $\left\{\begin{array}{l}y\le x\\x\le3\\y\ge0\end{array}\right.$
  \task $\left\{\begin{array}{l}x+y\ge2\\x\le4\\y\le3\end{array}\right.$
  \task $\left\{\begin{array}{l}2x+y\le6\\x+y\ge2\\x\ge0,\ y\ge0\end{array}\right.$
  \task $\left\{\begin{array}{l}x+y\ge3\\y\ge1\\x\ge0\end{array}\right.$
\end{tasks}\recorda{Cada vértice es el corte de dos rectas frontera: resuelve el sistema formado por sus ecuaciones.}

\ex{2} ¿Tiene solución el sistema $\left\{\begin{array}{l}x+2y\le2\\ x+2y\ge6\end{array}\right.$? Razónalo.

\ex{3} Escribe un sistema de inecuaciones cuya solución sea el rectángulo de vértices $(1,1)$, $(4,1)$, $(4,3)$ y $(1,3)$, con los bordes incluidos.


% =====================================================
\seccio{10}{Programación lineal}

\exf{1} Calcula el máximo de $F(x,y)=3x+2y$ en la región definida por $x+y\le5$, $x\le3$, $x\ge0$, $y\ge0$.

\ex{2} Calcula el mínimo de $F(x,y)=2x+3y$ con las restricciones $x+y\ge4$, $x\le5$, $y\le5$, $x\ge0$, $y\ge0$.

\ex{3} Maximiza $F(x,y)=x+2y$ con $x+2y\le8$, $x\le4$, $x\ge0$, $y\ge0$. ¿Qué observas?

\ex{2} Una pastelería hace pasteles de chocolate y de nata. Cada día puede hacer como máximo 20, y como máximo 12 de chocolate. Cada pastel de chocolate da 6 € de beneficio y cada pastel de nata, 4 €. ¿Cuántos tiene que hacer de cada tipo para obtener el beneficio máximo?

\ex{2} Un taller hace mesas y sillas. Cada mesa necesita 2 horas de carpintería y 1 de pintura, y cada silla, 1 hora de carpintería y 1 de pintura. Cada semana hay 40 horas de carpintería y 30 de pintura. El beneficio es de 50 € por mesa y de 30 € por silla. ¿Cuántas mesas y sillas tiene que hacer para ganar el máximo?

\ex{2} Un dietista combina dos alimentos, A y B. Cada ración de A tiene 1 unidad de hierro y 3 de calcio, y cada ración de B, 2 de hierro y 1 de calcio. Hay que tomar como mínimo 10 unidades de hierro y 15 de calcio. Si una ración de A cuesta 2 € y una de B, 3 €, ¿cuántas raciones de cada alimento hacen la dieta más barata?

\ex{3} Una persona quiere invertir como máximo 10\,000 € en dos fondos. El fondo A da un 4\,\% anual y el B, un 6\,\%. En el fondo B quiere poner como máximo 6\,000 € y en el A, como mínimo 2\,000 €. ¿Cómo tiene que repartir el dinero para obtener el máximo beneficio? (Trabaja en miles de euros.)

\ex{3} Un instituto organiza una excursión para 400 alumnos. La empresa de transporte tiene 8 autobuses de 40 plazas y 10 de 50 plazas, pero solo 9 conductores. Alquilar un autobús pequeño cuesta 600 € y uno grande, 800 €. ¿Cuántos autobuses de cada tipo hay que alquilar para que la excursión sea lo más barata posible?

\vfill
\begin{center}\color{muted}\small\faCheckCircle[regular]\ \ ¡Ánimo! Comprueba tus resultados con las herramientas interactivas de los apuntes de la unidad.\end{center}

\end{document}
